% !TeX program = lualatex
% ===================================================================
% mwe_acessivel.tex — teste mínimo do livros_crp_acessivel.sty
%
% Objetivo: verificar layout, fontes, page styles, capítulos,
% seções, partes e bookmarks antes de portar a seção 7 do pacote.
%
% Compilar com: lualatex mwe_acessivel.tex (duas vezes)
%
% Pré-requisitos no diretório:
%   - livros_crp_acessivel.sty
%   - Lora-Regular.ttf, Lora-Bold.ttf, Lora-Italic.ttf, Lora-BoldItalic.ttf
%   - NEWJUNE-*.OTF (todas as variantes usadas)
%   - Um arquivo de imagem qualquer para o teste de \includegraphics
%     (ajustar nome na linha indicada ou comentar o bloco)
% ===================================================================

\DocumentMetadata{
    lang        = pt-BR,
    pdfstandard = ua-2,
    pdfversion  = 2.0,
    testphase   = {phase-III,table,firstaid},
    tagging     = on,
}

\documentclass[11pt,a5paper,twoside,openright]{scrbook}

\usepackage{livros_crp_acessivel}

% Metadados do documento (vão para XMP via \DocumentMetadata + hyperref)
\hypersetup{
    pdftitle  = {MWE — Teste de acessibilidade},
    pdfauthor = {CRP-SP},
    pdflang   = {pt-BR},
}

\title{MWE — Teste de acessibilidade}
\author{CRP-SP}

\begin{document}

\frontmatter
\pagestyle{pretext}

% --- Folha de rosto simplificada ---
\begin{titlepage}
    \centering
    \vspace*{\fill}
    {\huge\bfseries\color{crp1} MWE — Teste de acessibilidade\par}
    \vspace{2\baselineskip}
    {\Large Verificação do pacote \texttt{livros\_crp\_acessivel}\par}
    \vspace{\fill}
    {\large CRP-SP\par}
    \vspace{2\baselineskip}
\end{titlepage}

% --- Sumário ---
\tableofcontents

\mainmatter
\pagestyle{scrheadings}

% ===================================================================
% PARTE I — teste de \part com número e título
% ===================================================================
\part{Primeira parte do teste}

% ===================================================================
% Capítulo 1 — texto corrido, teste de fontes e hifenização
% ===================================================================
\chapter{Texto corrido e tipografia}

Este capítulo existe para verificar se a fonte Lora está sendo
carregada corretamente como fonte principal, se o espaçamento
entre linhas corresponde ao \texttt{\textbackslash onehalfspacing},
e se a hifenização portuguesa está ativa por meio do pacote
\texttt{babel}.

O próximo parágrafo contém um texto longo o suficiente para que
o algoritmo de quebra de linhas do \LaTeX{} precise hifenizar
pelo menos uma palavra composta, como
\emph{extraordinariamente} ou \emph{responsabilidade}, e assim
demonstrar que a hifenização está se comportando conforme o
esperado no contexto do idioma selecionado.

Aqui testamos \textbf{negrito}, \emph{itálico},
\textsf{fonte sem serifa (NEWJUNE)}, e o peso customizado
{\Light peso light da NEWJUNE} seguido por
{\media peso medium da NEWJUNE} e por fim
{\fontebook peso book da NEWJUNE}.

% ===================================================================
% Seção — teste de hierarquia para tagging
% ===================================================================
\section{Seção de primeiro nível}

Este texto está dentro de uma seção. Em PDF tagueado, deve
aparecer como \texttt{<H2>} na árvore de estrutura (já que
\texttt{<H1>} corresponde ao capítulo).

\subsection{Subseção}

Texto da subseção. Deve virar \texttt{<H3>} nas tags.

\subsubsection{Sub-subseção}

Texto da sub-subseção. Deve virar \texttt{<H4>} nas tags.

\section{Outra seção}

Esta seção serve para verificar se há numeração visível (não
deve haver — \texttt{secnumdepth=-1}) e se a navegação por
bookmarks do \texttt{hyperref} lista corretamente capítulos
e seções.

% ===================================================================
% Capítulo 2 — elementos variados
% ===================================================================
\chapter{Elementos variados}

\section{Listas}

Lista com marcadores:
\begin{itemize}
    \item Primeiro item da lista.
    \item Segundo item, um pouco mais longo para testar quebra
          de linha dentro de item.
    \item Terceiro item.
\end{itemize}

Lista numerada:
\begin{enumerate}
    \item Primeiro.
    \item Segundo.
    \item Terceiro.
\end{enumerate}

\section{Tabela simples}

Uma tabela curta para verificar tagging automático via
\texttt{testphase=table}:

\begin{center}
\begin{tabular}{lcr}
    \toprule
    \textbf{Item} & \textbf{Quantidade} & \textbf{Valor} \\
    \midrule
    Livro         & 10                  & R\$ 250,00     \\
    Caderno       &  5                  & R\$  75,00     \\
    Caneta        & 20                  & R\$  40,00     \\
    \bottomrule
\end{tabular}
\end{center}

\section{Imagem com alt text}

% Comentar este bloco se não houver imagem disponível.
% Substituir "exemplo" pelo nome real de um arquivo de imagem.
%
% \begin{figure}[h]
%     \centering
%     \includegraphics[width=0.5\textwidth, alt={Descrição da imagem para leitor de tela}]{exemplo}
%     \caption{Legenda da figura de teste.}
% \end{figure}

Se houver uma imagem no bloco acima (descomentado), ela deve
aparecer na árvore de tags como \texttt{<Figure>} com o
atributo \texttt{/Alt} preenchido.

\section{Link e URL}

Um link externo para testar o \texttt{hyperref}:
\url{https://www.crpsp.org.br}

E um link com texto customizado:
\href{https://www.crpsp.org.br}{site do CRP-SP}.

% ===================================================================
% Capítulo 3 — página curta para testar paginação
% ===================================================================
\chapter{Página final}

Último capítulo do MWE. Serve para confirmar que a numeração
de páginas no cabeçalho aparece em crp1, que o cabeçalho usa
a fonte \texttt{\textbackslash Light} em tamanho \texttt{small},
e que a régua separadora do cabeçalho está na cor correta.

\end{document}
